\chapter{The Deterministic Standard-Basis Lift}\label{ch07:chap}

\section{Overview and evidential status}\label{ch07:sec:overview}

The preceding chapters built the typed finite-function category: named objects
with explicitly ordered carriers, total arrows with declared source and target,
partial type-checked composition, and the hard separation between endomaps
(which carry monoid powers) and semigroupoid paths (which deliberately do
not). This chapter develops the kernel's one piece of genuine linear algebra:
the \emph{deterministic standard-basis lift}, which assigns to every typed
arrow $f:X\to Y$ a linear map $A_f$ between finite-dimensional Hilbert spaces,
acting on basis vectors by
\[
A_f e_x = e_{f(x)} .
\]
The lift is deliberately elementary. It introduces no probability, no
dynamics, and no physical interpretation. Its purpose inside the kernel is
threefold. First, it provides an \emph{exactly analyzable} numerical shadow of
every arrow: the matrix of $A_f$ is a $0/1$ matrix with exactly one $1$ per
column, every spectral quantity of interest is an integer or the square root
of an integer, and the implementation in
\texttt{src/tks\_v8\_core/deterministic\_lift.py} computes all of them in exact
integer arithmetic, converting to floating point only at the single property
\texttt{operator\_norm}. Second, it supplies a family of \emph{permanent
counterexamples} that block two tempting but false general claims about
operator lifts of functions---claims that, if imported carelessly into
downstream CAS reasoning, would license contractivity and projection
arguments that are simply wrong. Third, its functoriality
($A_{\id_X}=I$, $A_{g\circ f}=A_g A_f$) gives an independent, exactly
checkable consistency test of the categorical composition of
Chapter~\ref{ch04:main}: the executable suite verifies the functor laws by
exact integer matrix equality on $1{,}678$ composition pairs.

The evidential status of every claim in this chapter is stated where the claim
is made, following the ledger discipline of the governing sources. In summary:

\begin{itemize}
\item The definition of the lift and its elementary structural facts
  (Section~\ref{ch07:sec:lift}) are definitions and immediate computations.
\item The exact operator-norm formula
  $\lVert A_f\rVert^2=\max_{y\in Y}|f^{-1}(y)|$
  (Theorem~\ref{ch07:thm:norm}, ledger record
  \textsc{thm-lift-norm-fiber-formula}) carries an \emph{informal proof},
  reproduced here in fully expanded form, together with bounded exhaustive
  structural tests. Its universal formalization in a proof assistant is
  \emph{open} (obligation PO-TKS8-CORE-05): the pinned Lean seed kernel-checks
  only the exact two-state expansion of
  Section~\ref{ch07:sec:counterexample}, not the general theorem.
\item The corollaries (trace, rank, isometry) inherit the same status; they
  are short consequences of the norm analysis and of the fiber decomposition.
\item The permanent counterexample (ledger record
  \textsc{thm-lift-occupation-counterexample}) is a finite, fully expanded
  computation; the Lean seed kernel-checks this exact two-state expansion and
  its noncontractivity.
\item Functoriality (Theorem~\ref{ch07:thm:functor}, obligation
  PO-TKS8-CORE-11) carries \emph{partial formal verification}: the Lean seed
  kernel-checks the identity lift entrywise and on basis vectors, typed
  composition on basis vectors, and relational-product support composition;
  the finite-sum numerical matrix bridge and arbitrary-vector linearity
  remain open.
\item No stochastic, quantum-channel, contractivity, dynamical, or empirical
  claim is made anywhere in this chapter; Section~\ref{ch07:sec:notclaims}
  is devoted to stating these non-claims precisely.
\end{itemize}

This chapter, like the whole monograph, is a Ring~4 draft artifact: nothing
here is peer-reviewed, nothing is licensed for semantic transfer onto
canonical TKS/CAS constructs, and the drafting of the chapter promotes no
claim beyond what the underlying Gate~1 sources already record.

\section{Finite-dimensional preliminaries}\label{ch07:sec:prelim}

We collect the small amount of linear algebra the chapter needs, in the exact
form in which the proofs use it. Everything in this section is standard; the
canonical reference for the spectral background is Horn and
Johnson~\cite{horn}. We restate the facts with proofs where the proofs are
short enough to keep the chapter self-contained, because the norm theorem
below is a registered ledger item and we want its complete logical
dependency chain visible on the page.

\begin{notation}[Spaces and bases]\label{ch07:not:spaces}
For a finite object $X$ with ordered carrier $(x_0,\dots,x_{|X|-1})$, let
$H_X$ denote the complex inner-product space $\mathbb{C}^{|X|}$ with the
standard Hermitian inner product
$\langle u,v\rangle=\sum_{i} \overline{u_i}\,v_i$
and with orthonormal basis $(e_{x_0},\dots,e_{x_{|X|-1}})$ indexed by the
ordered carrier; $e_{x_i}$ is the $i$-th standard coordinate vector. The
ordered-input discipline of Chapter 3 is what makes this
indexing well defined: because the carrier is an explicitly ordered sequence,
the correspondence between carrier elements and basis indices is a fixed,
replayable bijection, not an artifact of interpreter hash order. When $X$ is
empty, $H_X$ is the zero space.
\end{notation}

\begin{definition}[Adjoint]\label{ch07:def:adjoint}
For a linear map $A:H_X\to H_Y$, the \emph{adjoint} $A^*:H_Y\to H_X$ is the
unique linear map satisfying
$\langle A^* w, u\rangle = \langle w, A u\rangle$
for all $u\in H_X$, $w\in H_Y$. In matrix terms, $A^*$ is the conjugate
transpose of $A$; for real (in particular $0/1$) matrices, the plain
transpose.
\end{definition}

\begin{definition}[Operator norm]\label{ch07:def:opnorm}
The \emph{operator norm} of $A:H_X\to H_Y$ is
\[
\lVert A\rVert
=\sup\{\lVert A u\rVert : u\in H_X,\ \lVert u\rVert\le 1\}
=\max\{\lVert A u\rVert : \lVert u\rVert = 1\}
\]
when $H_X$ is nonzero (the supremum is attained by compactness of the unit
sphere in finite dimension), and $\lVert A\rVert=0$ when $H_X$ is the zero
space. The second form is the one we use; the empty-domain convention matches
the implementation, whose \texttt{max\_fiber\_size} takes
\texttt{max(\ldots, default=0)}.
\end{definition}

\begin{lemma}[Norm via the Gram operator]\label{ch07:lem:gramreduction}
For any linear $A:H_X\to H_Y$ with $H_X$ nonzero,
$\lVert A\rVert^2=\lambda_{\max}(A^*A)$, the largest eigenvalue of the
positive semidefinite Hermitian operator $A^*A$.
\end{lemma}

\begin{proof}
For every unit vector $u$,
$\lVert Au\rVert^2=\langle Au,Au\rangle=\langle A^*Au,u\rangle$, so
$\lVert A\rVert^2=\max_{\lVert u\rVert=1}\langle A^*Au,u\rangle$. The
operator $G=A^*A$ is Hermitian ($G^*=(A^*A)^*=A^*A=G$) and positive
semidefinite ($\langle Gu,u\rangle=\lVert Au\rVert^2\ge 0$). By the spectral
theorem for Hermitian matrices \cite{horn}, $H_X$ admits an orthonormal basis
$(v_1,\dots,v_n)$ of eigenvectors of $G$ with real eigenvalues
$\lambda_1\ge\dots\ge\lambda_n\ge 0$. Writing a unit vector as
$u=\sum_k c_k v_k$ with $\sum_k |c_k|^2=1$,
\[
\langle Gu,u\rangle=\sum_k \lambda_k |c_k|^2
\le \lambda_1 \sum_k |c_k|^2 = \lambda_1,
\]
with equality at $u=v_1$. Hence the maximum of the quadratic form on the unit
sphere is exactly $\lambda_1=\lambda_{\max}(G)$.
\end{proof}

\begin{lemma}[Spectrum is a similarity invariant]\label{ch07:lem:similarity}
If $P$ is a permutation matrix and $G$ any square matrix of matching size,
then $PGP^{-1}$ has the same characteristic polynomial, hence the same
eigenvalues with the same algebraic multiplicities, as $G$. Moreover, when
$G$ is Hermitian, so is $PGP^{-1}$, because a real permutation matrix is
unitary ($P^{-1}=P^{T}=P^{*}$).
\end{lemma}

\begin{proof}
$\det(tI-PGP^{-1})=\det\bigl(P(tI-G)P^{-1}\bigr)
=\det(P)\det(tI-G)\det(P^{-1})=\det(tI-G)$. For the Hermitian statement,
$(PGP^{*})^{*}=PG^{*}P^{*}=PGP^{*}$.
\end{proof}

\begin{lemma}[Spectra of direct sums]\label{ch07:lem:directsum}
If $G=\bigoplus_j G_j$ is block diagonal, the spectrum of $G$ is the
multiset union of the spectra of the blocks $G_j$.
\end{lemma}

\begin{proof}
The characteristic polynomial of a block-diagonal matrix factors as the
product of the characteristic polynomials of the blocks:
$\det(tI-G)=\prod_j\det(tI_j-G_j)$. Eigenvalues and algebraic multiplicities
therefore add up across blocks.
\end{proof}

These three lemmas are the entire spectral toolkit the norm theorem needs.
One further standard notion appears in the counterexample analysis:

\begin{definition}[Projector, partial isometry]\label{ch07:def:projpi}
A linear operator $Q$ on a Hilbert space is an \emph{orthogonal projector}
when $Q^2=Q$ and $Q^*=Q$. A linear map $A:H_X\to H_Y$ is a \emph{partial
isometry} when $A A^* A=A$; equivalently, when $A^*A$ is an orthogonal
projector, equivalently when $A$ restricts to an isometry on the orthogonal
complement of its kernel \cite{horn}. An \emph{isometry} is a map with
$A^*A=I$, i.e.\ $\lVert Au\rVert=\lVert u\rVert$ for all $u$.
\end{definition}

Note the direction of the traps already visible here: idempotence ($Q^2=Q$)
alone does \emph{not} make an orthogonal projector---self-adjointness is a
separate condition---and having $0/1$ entries does not make a partial
isometry. The counterexample of Section~\ref{ch07:sec:counterexample} lives
exactly in this gap.

\section{The lift: definition, matrix form, and running examples}\label{ch07:sec:lift}

\subsection{Definition}

\begin{definition}[Deterministic standard-basis lift]\label{ch07:def:lift}
Let $f:X\to Y$ be an arrow of the typed finite-function category: a total map
sending each standard index of the ordered carrier of $X$ to a standard index
of the ordered carrier of $Y$. The \emph{deterministic standard-basis lift}
of $f$ is the linear map
\[
A_f:H_X\longrightarrow H_Y,\qquad A_f e_x = e_{f(x)}
\quad\text{for every carrier element } x \text{ of } X,
\]
extended to all of $H_X$ by linearity:
$A_f\bigl(\textstyle\sum_x c_x e_x\bigr)=\sum_x c_x\, e_{f(x)}$.
\end{definition}

The definition is well posed for exactly the reasons the kernel's object
discipline was designed to secure. Because $f$ is \emph{total}, every basis
vector of $H_X$ has an image; because $f$ lands in the declared target's
standard indices, that image is a basis vector of $H_Y$; because carriers are
\emph{explicitly ordered}, the phrase ``the $i$-th basis vector'' has a fixed
meaning independent of interpreter state; and because a linear map is
determined by its values on a basis, $A_f$ is unique. Nothing else about $f$
is used: in particular no injectivity, surjectivity, or endomorphism
assumption. Every typed arrow lifts.

\begin{proposition}[Matrix form]\label{ch07:prop:matrix}
With respect to the standard bases, the matrix of $A_f$ is the
$|Y|\times|X|$ matrix with entries
\[
(A_f)_{y x}=
\begin{cases}
1 & \text{if } f(x)=y,\\
0 & \text{otherwise.}
\end{cases}
\]
Each column contains exactly one $1$ (in the row indexed by the image of that
column's carrier element) and each column therefore sums to $1$. Conversely,
every $|Y|\times|X|$ matrix of $0$s and $1$s with exactly one $1$ per column
is the lift matrix of exactly one arrow $f:X\to Y$.
\end{proposition}

\begin{proof}
Column $x$ of the matrix of $A_f$ is the coordinate vector of $A_f e_x =
e_{f(x)}$, which has a single $1$ in row $f(x)$. For the converse, given such
a matrix $M$, define $f(x)$ to be the unique row index with $M_{f(x),x}=1$;
this reads off a total map, and its lift matrix is $M$ by the first part. The
correspondence $f\mapsto A_f$ is therefore a bijection between arrows
$X\to Y$ and one-$1$-per-column $0/1$ matrices of shape $|Y|\times|X|$; in
particular the lift is \emph{faithful}: distinct arrows with the same source
and target have distinct lifts.
\end{proof}

The implementation realizes Definition~\ref{ch07:def:lift} as the frozen
dataclass \texttt{DeterministicLift} wrapping a \texttt{FiniteArrow}; its
method \texttt{matrix()} returns the tuple-of-tuples integer matrix of
Proposition~\ref{ch07:prop:matrix}, its method \texttt{apply(vector)}
computes $A_f$ on an explicit ordered coordinate sequence by accumulating
$c_x$ into output slot $f(x)$, and \texttt{adjoint\_apply(vector)} computes
$A_f^*$ on target coordinates by reading coordinate $f(x)$ back into slot
$x$. Both methods enforce the ordered-input discipline: a set, mapping,
string, or one-shot iterable in coordinate position raises
\texttt{ValidationError}, for the same reason carriers do---coordinate order
is basis order, and basis order is meaning.

\begin{remark}[The adjoint as pullback]\label{ch07:rem:adjointpullback}
The formula implemented by \texttt{adjoint\_apply} deserves a second look,
because it identifies the adjoint with something combinatorially natural.
On basis vectors,
\[
A_f^* e_y=\sum_{x\,:\,f(x)=y} e_x ,
\]
the indicator sum of the fiber $f^{-1}(y)$. Indeed
$\langle A_f^* e_y, e_x\rangle = \langle e_y, A_f e_x\rangle
=\langle e_y, e_{f(x)}\rangle$, which is $1$ iff $f(x)=y$. On general
coordinates, $(A_f^* w)_x = w_{f(x)}$: the adjoint is precomposition with
$f$, i.e.\ the \emph{pullback} of coordinate functions, while $A_f$ itself is
the \emph{pushforward} of coordinates. The pushforward merges coordinates
that share a fiber (adding them); the pullback duplicates a coordinate across
its fiber. All four derived operators of Section~\ref{ch07:sec:derived} are
compositions of one merge and one duplication in the two possible orders.
\end{remark}

\begin{example}[Coordinates through the lift, numerically]\label{ch07:ex:coords}
To fix the mechanics of \texttt{apply} and \texttt{adjoint\_apply}, push one
concrete vector each way through the arrow $f:X_4\to Y_3$ of the running
quartet below ($f(p)=f(q)=a$, $f(r)=f(s)=b$). Take source coordinates
$(2,\,3,\,5,\,7)$ in carrier order $(p,q,r,s)$. Then
\[
A_f(2,3,5,7)^{T}=(2+3,\ 5+7,\ 0)^{T}=(5,12,0)^{T}:
\]
the coordinates of $p$ and $q$ \emph{merge additively} into slot $a$, those
of $r$ and $s$ into slot $b$, and slot $c$ --- an empty fiber --- receives
nothing. In the other direction, take target coordinates $(10,20,30)$ in
carrier order $(a,b,c)$. Then
\[
A_f^{*}(10,20,30)^{T}=(10,\,10,\,20,\,20)^{T}:
\]
each source slot \emph{reads back} the coordinate of its image, so the value
at $a$ is duplicated across the fiber $\{p,q\}$ and the value at $c$ is
simply never consulted. Merge-by-addition and duplicate-by-fiber are the two
primitive motions; every operator in this chapter is a composite of them,
and every collision cost the norm theorem quantifies is visible already in
the first computation, where the input mass $2+3$ arrived at $a$
constructively. (Replacing $(2,3)$ by $(1,1)/\sqrt2$ exhibits the extremal
vector of the norm theorem for this fiber.)
\end{example}

\subsection{Running examples}\label{ch07:sec:running}

Four arrows accompany us through the whole chapter. They are chosen to
exhaust the qualitative possibilities: non-injective and non-surjective,
injective and non-surjective, bijective, and maximally collapsing.

\begin{example}[The running quartet]\label{ch07:ex:quartet}
Fix three objects with the indicated ordered carriers:
\[
X_4=(\texttt{State4},\ (p,q,r,s)),\qquad
X_2=(\texttt{Pair},\ (u,v)),\qquad
Y_3=(\texttt{Obs3},\ (a,b,c)).
\]
Define four arrows:
\begin{itemize}
\item[(R1)] $f:X_4\to Y_3$ with $f(p)=a$, $f(q)=a$, $f(r)=b$, $f(s)=b$.
  Fibers: $f^{-1}(a)=\{p,q\}$, $f^{-1}(b)=\{r,s\}$, $f^{-1}(c)=\emptyset$.
  Neither injective nor surjective.
\item[(R2)] $g:X_2\to Y_3$ with $g(u)=c$, $g(v)=a$. Injective, not
  surjective.
\item[(R3)] $\pi:X_3\to X_3$ on a three-element object
  $X_3=(\texttt{Cycle3},(0,1,2))$ with $\pi(0)=1$, $\pi(1)=2$, $\pi(2)=0$.
  Bijective (a $3$-cycle).
\item[(R4)] $c_0:X_2'\to X_2'$ the constant map on a two-element object
  $X_2'=(\texttt{Bit},(0,1))$ with $c_0(0)=c_0(1)=0$. Maximally collapsing;
  this is the permanent counterexample of
  Section~\ref{ch07:sec:counterexample}.
\end{itemize}
Their lift matrices, rows indexed by target carrier order and columns by
source carrier order, are
\[
A_f=\begin{pmatrix}1&1&0&0\\0&0&1&1\\0&0&0&0\end{pmatrix},\qquad
A_g=\begin{pmatrix}0&1\\0&0\\1&0\end{pmatrix},\qquad
A_\pi=\begin{pmatrix}0&0&1\\1&0&0\\0&1&0\end{pmatrix},\qquad
A_{c_0}=\begin{pmatrix}1&1\\0&0\end{pmatrix}.
\]
Each matrix has exactly one $1$ per column; the reader should check one entry
of each against the defining rule $(A)_{yx}=1$ iff the arrow sends $x$ to
$y$. Note how the empty fiber of $f$ at $c$ appears as an all-zero \emph{row}
(rows are free to be zero; columns are not), and how the bijectivity of
$\pi$ appears as ``exactly one $1$ per row'' \emph{in addition to} one per
column---i.e.\ a permutation matrix.
\end{example}

\begin{remark}[What the lift forgets and what it keeps]\label{ch07:rem:forget}
The lift keeps everything: by
Proposition~\ref{ch07:prop:matrix} it is a bijection onto its matrix class,
so $f$ is recoverable from $A_f$ by reading each column. What changes is the
\emph{ambient theory}. As a function, $f$ supports composition, powers (if
endo), and fiber counting. As a matrix, $A_f$ additionally supports addition,
scalar action, adjoints, spectra, and norms---structure that has no preimage
in the category of functions. The lift is thus an embedding of a small
combinatorial world into a larger linear one, and every use the kernel makes
of it takes the form: compute a linear invariant, then translate the answer
back into fiber combinatorics. The norm theorem below is the paradigm case.
\end{remark}

\subsection{Degenerate shapes: empty carriers}\label{ch07:sec:degenerate}

The kernel admits empty carriers, and the lift theory must therefore say
something coherent about them rather than excluding them by fiat; the
executable evidence deliberately includes empty carriers in its
$1{,}678$ functor-law pairs, so the conventions here are load-bearing.

\emph{Empty source, nonempty target.} If $|X|=0$ there is exactly one arrow
$f:X\to Y$ for each $Y$ (the empty map: totality quantifies over an empty
carrier and is vacuously satisfied). Its lift is the unique linear map from
the zero space to $H_Y$, namely the zero map; its matrix has shape
$|Y|\times 0$ --- $|Y|$ rows and no columns --- and the
one-$1$-per-column rule is vacuously true. Every fiber is empty, the fiber
maximum is $0$ under the stated convention, and
$\lVert A_f\rVert=0$ consistently with Definition~\ref{ch07:def:opnorm}.
The input Gram is the $0\times 0$ matrix; the output Gram and the support
projector are the $|Y|\times|Y|$ zero matrix; the trace-Gram is
$0=|X|$; the rank is $0=|\im f|$. Every corollary of this chapter holds
degenerately, which is exactly why the convention was chosen as it was.

\emph{Empty target.} An arrow into an empty object exists only when the
source is also empty: a total map must send each of the $|X|$ carrier
elements somewhere, and there is nowhere. The typed constructor enforces
this (there is no index to map to), so the case $|Y|=0<|X|$ simply does not
arise, and no convention is needed for it. The one arrow
$\emptyset\to\emptyset$ lifts to the identity of the zero space, which is
simultaneously the zero map --- the unique point where those two agree.

\emph{Is the empty-source lift injective?} The implementation answers yes:
\texttt{is\_injective} checks that all fiber sizes are at most $1$, which
holds vacuously, and correspondingly \texttt{is\_isometry} reports true.
This is the mathematically standard reading (the empty map is injective; a
zero-space map preserves all norms there are), but it produces the one
place where a slogan of this chapter needs a footnote: ``isometries have
norm $1$'' holds only for nonempty sources, since the empty lift is an
isometry of norm $0$. Corollary~\ref{ch07:cor:isometry} is stated with the
nonemptiness hypothesis for precisely this reason. Degenerate cases do not
threaten theorems; unexamined degenerate cases threaten slogans.

\section{Three derived operators, kept distinct}\label{ch07:sec:derived}

From $A_f$ the implementation derives three operators. Informal usage tends
to conflate them---all three are ``the projection-like thing associated to
$f$''---and part of the kernel's discipline is to keep them apart as three
distinct methods with three distinct semantics. We define each, compute all
of them on the running examples, and record exactly when any two coincide.

\subsection{The input-collision Gram matrix}

\begin{definition}[Input Gram]\label{ch07:def:inputgram}
The \emph{input-collision Gram matrix} of $f$ is
$G_f = A_f^* A_f : H_X\to H_X$.
\end{definition}

\begin{proposition}[Gram entries]\label{ch07:prop:gramentries}
$(G_f)_{x x'} = 1$ if $f(x)=f(x')$ and $0$ otherwise. In particular every
diagonal entry is $1$, and $G_f$ is the block-indicator matrix of the
partition of the source carrier into fibers of $f$.
\end{proposition}

\begin{proof}
$(G_f)_{xx'}=\langle e_x, A_f^*A_f e_{x'}\rangle
=\langle A_f e_x, A_f e_{x'}\rangle
=\langle e_{f(x)}, e_{f(x')}\rangle$,
which is $1$ iff $f(x)=f(x')$ by orthonormality of the target basis.
Diagonal entries take $x=x'$. The relation ``$f(x)=f(x')$'' is an
equivalence relation on the source carrier whose classes are exactly the
nonempty fibers, so $G_f$ is the $0/1$ indicator of that partition.
\end{proof}

The name is chosen with care: $G_f$ records \emph{input collisions}. It is an
operator on the \emph{source} space, it sees only which inputs are merged,
and it is insensitive to which target element they are merged into and to
whether other target elements are hit at all. The implementation method is
\texttt{input\_gram\_matrix()}.

\begin{example}[Input Grams of the quartet]\label{ch07:ex:gramquartet}
\[
G_f=\begin{pmatrix}1&1&0&0\\1&1&0&0\\0&0&1&1\\0&0&1&1\end{pmatrix},\qquad
G_g=\begin{pmatrix}1&0\\0&1\end{pmatrix},\qquad
G_\pi=\begin{pmatrix}1&0&0\\0&1&0\\0&0&1\end{pmatrix},\qquad
G_{c_0}=\begin{pmatrix}1&1\\1&1\end{pmatrix}.
\]
For $f$, the two $2\times 2$ all-ones blocks are the fibers $\{p,q\}$ and
$\{r,s\}$; the carrier order happens to list each fiber contiguously, so the
block structure is visible without reordering (in general it is not; the
norm proof below handles the general case by an explicit permutation). For
injective $g$ and bijective $\pi$ the Gram is the identity: no collisions.
For the constant map the Gram is the all-ones matrix $J_2$: total collision.
\end{example}

\subsection{The output Gram: fiber sizes on the diagonal}

\begin{definition}[Output Gram]\label{ch07:def:outputgram}
The \emph{output Gram matrix} of $f$ is
$\widehat{G}_f = A_f A_f^* : H_Y\to H_Y$.
\end{definition}

\begin{proposition}[Output Gram entries]\label{ch07:prop:outputgram}
$\widehat{G}_f$ is diagonal, with
$(\widehat{G}_f)_{yy}=|f^{-1}(y)|$, the size of the fiber over $y$.
\end{proposition}

\begin{proof}
$(\widehat{G}_f)_{y y'}=\langle e_y, A_fA_f^* e_{y'}\rangle
=\langle A_f^* e_y, A_f^* e_{y'}\rangle$. By
Remark~\ref{ch07:rem:adjointpullback}, $A_f^*e_y$ is the indicator vector of
the fiber $f^{-1}(y)$. Two distinct fibers are disjoint, so their indicator
vectors are orthogonal, giving $0$ off the diagonal; and the inner product of
a fiber indicator with itself is the fiber's cardinality.
\end{proof}

So $\widehat{G}_f$ lives on the \emph{target} space and counts \emph{how
many} inputs land on each output. Its diagonal is exactly the tuple computed
by the implementation property \texttt{fiber\_sizes} (an integer tuple of
length $|Y|$, obtained by a single counting pass over the mapping), and the
method \texttt{output\_gram\_matrix()} places those counts on a diagonal.

\begin{example}[Output Grams of the quartet]\label{ch07:ex:outputquartet}
\[
\widehat{G}_f=\begin{pmatrix}2&0&0\\0&2&0\\0&0&0\end{pmatrix},\qquad
\widehat{G}_g=\begin{pmatrix}1&0&0\\0&0&0\\0&0&1\end{pmatrix},\qquad
\widehat{G}_\pi=I_3,\qquad
\widehat{G}_{c_0}=\begin{pmatrix}2&0\\0&0\end{pmatrix}.
\]
The diagonals read off the fiber-size tuples $(2,2,0)$, $(1,0,1)$,
$(1,1,1)$, and $(2,0)$ respectively.
\end{example}

\subsection{The codomain support projector}

\begin{definition}[Support projector]\label{ch07:def:support}
The \emph{codomain support projector} of $f$ is the diagonal $0/1$ matrix
$S_f:H_Y\to H_Y$ with
$(S_f)_{yy}=1$ if $f^{-1}(y)\neq\emptyset$ and $0$ otherwise.
\end{definition}

\begin{proposition}\label{ch07:prop:supportisproj}
$S_f$ is an orthogonal projector, and its range is
$\im A_f=\operatorname{span}\{e_y : y\in\im f\}$.
\end{proposition}

\begin{proof}
A real diagonal matrix is self-adjoint, and a diagonal $0/1$ matrix squares
to itself, so $S_f$ is an orthogonal projector. Its range is spanned by the
basis vectors at which its diagonal is $1$, i.e.\ $\{e_y: y\in\im f\}$. The
columns of $A_f$ are exactly the vectors $e_{f(x)}$, so $\im A_f$ is the span
of $\{e_y : y \in \im f\}$ as well.
\end{proof}

The implementation method is \texttt{output\_support\_projector()}, computed
by thresholding the fiber-size tuple at zero. Its docstring calls it ``the
orthogonal projector onto $\im(A_f)$ in the codomain,'' which
Proposition~\ref{ch07:prop:supportisproj} verifies.

\begin{example}[Support projectors of the quartet]\label{ch07:ex:supportquartet}
\[
S_f=\begin{pmatrix}1&0&0\\0&1&0\\0&0&0\end{pmatrix},\qquad
S_g=\begin{pmatrix}1&0&0\\0&0&0\\0&0&1\end{pmatrix},\qquad
S_\pi=I_3,\qquad
S_{c_0}=\begin{pmatrix}1&0\\0&0\end{pmatrix}.
\]
\end{example}

\subsection{When the operators coincide, and why the distinction is load-bearing}

Comparing Examples~\ref{ch07:ex:outputquartet}
and~\ref{ch07:ex:supportquartet}: for $g$ and $\pi$ the output Gram and the
support projector are \emph{equal}; for $f$ and $c_0$ they differ (the
former has a $2$ where the latter has a $1$). This is the general pattern:

\begin{proposition}[Coincidence criterion]\label{ch07:prop:coincide}
$\widehat{G}_f=S_f$ if and only if every fiber of $f$ has size at most $1$,
i.e.\ iff $f$ is injective. Likewise $G_f=I_{H_X}$ iff $f$ is injective, and
$G_f$ is an orthogonal projector iff its all-ones fiber blocks all have size
$1$, again iff $f$ is injective. For non-injective $f$, all three derived
operators are pairwise distinct objects: $G_f$ is generally non-diagonal;
$\widehat{G}_f$ is diagonal but not $0/1$; $S_f$ is diagonal and $0/1$.
\end{proposition}

\begin{proof}
$\widehat{G}_f$ and $S_f$ are both diagonal; they agree exactly when
$|f^{-1}(y)|=\mathbf{1}[f^{-1}(y)\neq\emptyset]$ for every $y$, i.e.\ when no
fiber has two or more elements. $G_f=I$ iff there are no off-diagonal
collisions, i.e.\ no two distinct inputs share a fiber. For the projector
claim about $G_f$: $G_f$ is always self-adjoint (Proposition
\ref{ch07:prop:gramentries} makes it a real symmetric matrix), and by the
block description $G_f^2$ has $(x,x')$ entry equal to the size of the common
fiber when $f(x)=f(x')$ and $0$ otherwise (this is a special case of the
computation $J_m^2=mJ_m$ proved in Lemma~\ref{ch07:lem:allones} below); so
$G_f^2=G_f$ iff every nonempty fiber has size $1$.
\end{proof}

\begin{boundary}[Why three methods, not one]\label{ch07:bd:threemethods}
The conflation risk is concrete. In informal operator talk, all three of
$G_f$, $\widehat{G}_f$, $S_f$ get called ``the projector associated with
$f$,'' and for injective $f$ the talk is harmless because
$G_f=I$ and $\widehat{G}_f=S_f$. The moment a single collision appears, the
three objects diverge, only $S_f$ remains a projector, and any argument that
silently substitutes one for another acquires a factor equal to a fiber size
--- which is precisely the quantity the norm theorem says governs the
operator norm. The implementation therefore exposes them as three separately
named methods (\texttt{input\_gram\_matrix}, \texttt{output\_gram\_matrix},
\texttt{output\_support\_projector}) and the test suite pins each to its own
closed form. No claim of interchangeability is made or licensed anywhere in
the kernel.
\end{boundary}

\section{The exact norm theorem}\label{ch07:sec:norm}

We now prove the chapter's central result. Ledger status, restated for
discipline: \textsc{thm-lift-norm-fiber-formula}; informal proof present
(this section \emph{is} that proof, fully expanded); bounded exhaustive
structural tests in the executable suite; universal formalization
\emph{open} --- the pinned Lean seed kernel-checks only the two-state
expansion of Section~\ref{ch07:sec:counterexample}, not this theorem.

\begin{theorem}[Exact operator norm of a deterministic lift]\label{ch07:thm:norm}
Let $f:X\to Y$ be a typed arrow with $X$ nonempty. Then
\[
\lVert A_f\rVert^2=\max_{y\in Y}\bigl|f^{-1}(y)\bigr|,
\qquad\text{hence}\qquad
\lVert A_f\rVert=\sqrt{\max_{y\in Y}\bigl|f^{-1}(y)\bigr|}.
\]
For empty $X$ the norm is $0$ by the convention of
Definition~\ref{ch07:def:opnorm}, matching the formula if the empty maximum
is read as $0$.
\end{theorem}

The proof proceeds in four steps, each isolated as a lemma so that its
logical role is visible: (1) reduce the norm to the spectrum of the input
Gram; (2) reorder the source basis by a permutation similarity that groups
fibers into contiguous blocks; (3) analyze the all-ones block $J_m$ from
scratch; (4) assemble the spectrum and take the maximum.

\subsection{Step 1: reduction to the Gram spectrum}

By Lemma~\ref{ch07:lem:gramreduction},
$\lVert A_f\rVert^2=\lambda_{\max}(G_f)$ where $G_f=A_f^*A_f$ is the
input-collision Gram matrix, whose entries
Proposition~\ref{ch07:prop:gramentries} computed exactly:
$(G_f)_{xx'}=\mathbf{1}[f(x)=f(x')]$. The problem is now purely spectral:
find the largest eigenvalue of the $0/1$ indicator matrix of the fiber
partition.

\subsection{Step 2: permutation similarity groups the fibers}

\begin{lemma}[Fiber-block form]\label{ch07:lem:fiberblock}
Let the nonempty fibers of $f$ have sizes $m_1,\dots,m_r$ (so
$r=|\im f|$ and $m_1+\dots+m_r=|X|$). There is a permutation matrix $P$ on
$H_X$ such that
\[
P\,G_f\,P^{*}=J_{m_1}\oplus J_{m_2}\oplus\cdots\oplus J_{m_r},
\]
a block-diagonal direct sum of all-ones blocks, where $J_m$ denotes the
$m\times m$ matrix every entry of which is $1$.
\end{lemma}

\begin{proof}
Choose any enumeration of the source carrier that lists the elements of the
first nonempty fiber first (in any order), then the second fiber, and so on;
formally, a bijection $\sigma$ from $\{0,\dots,|X|-1\}$ to itself such that
the $\sigma$-preimages of each fiber form a contiguous interval of indices.
Let $P$ be the permutation matrix of $\sigma$, i.e.\ $Pe_{x_i}=e_{x_{\sigma(i)}}$
reindexed appropriately; $P$ is real orthogonal, so $P^{*}=P^{-1}$. The
$(i,j)$ entry of $P G_f P^{*}$ is the $(\sigma^{-1}(i),\sigma^{-1}(j))$ entry
of $G_f$, which is $1$ iff the carrier elements at positions
$\sigma^{-1}(i)$ and $\sigma^{-1}(j)$ share a fiber. Under the chosen
enumeration, positions share a fiber iff they lie in the same contiguous
interval. Hence the conjugated matrix is $1$ exactly on the $r$ diagonal
square blocks whose sizes are the fiber sizes, and $0$ elsewhere: the stated
direct sum.
\end{proof}

By Lemma~\ref{ch07:lem:similarity}, this conjugation changes neither the
eigenvalues nor their multiplicities, so we may compute the spectrum of the
direct sum instead of the spectrum of $G_f$.

It is worth pausing on why an explicit permutation is needed at all, when
Example~\ref{ch07:ex:gramquartet} already displayed a block-diagonal Gram.
There, the declared carrier order happened to list each fiber contiguously.
The kernel guarantees no such thing in general: carrier order is declared by
the caller for identity and replay reasons (Chapter 3), not
for the convenience of any particular arrow, and a single object is typically
the source of many arrows whose fiber partitions are mutually incompatible
with any one ordering. The permutation-similarity step is the precise
bookkeeping that makes the ``group the fibers'' intuition into an argument
valid for every arrow simultaneously, at the cost of one application of
Lemma~\ref{ch07:lem:similarity}.

\subsection{Step 3: the all-ones block, from scratch}

\begin{lemma}[Spectral analysis of $J_m$]\label{ch07:lem:allones}
For $m\ge 1$, the $m\times m$ all-ones matrix $J_m$ satisfies:
\begin{enumerate}
\item $J_m=\mathbf{1}\mathbf{1}^{*}$ where $\mathbf{1}=(1,\dots,1)^{T}$;
  consequently $\rk J_m=1$.
\item $J_m^2=mJ_m$.
\item The spectrum of $J_m$ is: eigenvalue $m$ with multiplicity $1$
  (eigenvector $\mathbf{1}$), and eigenvalue $0$ with multiplicity $m-1$
  (eigenspace the hyperplane $\{v:\sum_i v_i=0\}$).
\end{enumerate}
\end{lemma}

\begin{proof}
(1) The outer product $\mathbf{1}\mathbf{1}^{*}$ has $(i,j)$ entry
$1\cdot\overline{1}=1$, which is $J_m$. Every column of $J_m$ equals
$\mathbf{1}$, so the column space is the line spanned by $\mathbf{1}$ and the
rank is exactly $1$ (it is not $0$ since $J_m\neq 0$ for $m\geq 1$).

(2) Using (1) and associativity,
$J_m^2=\mathbf{1}(\mathbf{1}^{*}\mathbf{1})\mathbf{1}^{*}
=\mathbf{1}\,(m)\,\mathbf{1}^{*}=mJ_m$,
since $\mathbf{1}^{*}\mathbf{1}=\sum_{i=1}^m 1=m$. (Entrywise, the $(i,j)$
entry of $J_m^2$ is $\sum_{k=1}^m 1\cdot 1=m$.)

(3) First, $J_m\mathbf{1}$ has every entry $\sum_j 1=m$, so
$J_m\mathbf{1}=m\mathbf{1}$: the all-ones vector is an eigenvector with
eigenvalue $m$. Second, if $\sum_i v_i=0$ then every entry of $J_m v$ is
$\sum_i v_i=0$, so the sum-zero hyperplane, which has dimension $m-1$, lies
in the kernel: eigenvalue $0$ with geometric multiplicity at least $m-1$.
These two eigenspaces are orthogonal
($\langle\mathbf{1},v\rangle=\sum_i \overline{1}\, v_i=0$ for sum-zero $v$)
and their dimensions add to $m$, so together they exhaust $\mathbb{C}^m$;
since $J_m$ is Hermitian, algebraic and geometric multiplicities agree, and
the spectrum is exactly $\{m\}\cup\{0\text{ with multiplicity } m-1\}$.
Alternatively and purely algebraically: by (2), every eigenvalue $\lambda$
satisfies $\lambda^2=m\lambda$, so $\lambda\in\{0,m\}$; the trace of $J_m$ is
$m$, and the trace is the eigenvalue sum, so $m$ occurs exactly once and $0$
fills the remaining $m-1$ slots.
\end{proof}

Note the two independent routes to the multiplicity count in part (3): the
geometric route (explicit eigenvectors, orthogonal complement) and the
algebraic route (minimal-polynomial constraint from $J_m^2=mJ_m$ plus a trace
count). The redundancy is intentional. The lemma is the computational heart
of a registered ledger theorem, and the monograph's policy for such
load-bearing steps is to give two proofs whenever two short proofs exist.

\subsection{Step 4: assembling the spectrum}

\begin{proof}[Proof of Theorem~\textup{\ref{ch07:thm:norm}}]
Let the nonempty fibers of $f$ have sizes $m_1,\dots,m_r$. By
Lemma~\ref{ch07:lem:gramreduction},
$\lVert A_f\rVert^2=\lambda_{\max}(G_f)$. By
Lemma~\ref{ch07:lem:fiberblock} and Lemma~\ref{ch07:lem:similarity}, $G_f$
has the same spectrum as $J_{m_1}\oplus\cdots\oplus J_{m_r}$. By
Lemma~\ref{ch07:lem:directsum} and Lemma~\ref{ch07:lem:allones}, that
spectrum is the multiset
\[
\{\,m_1,\dots,m_r\,\}\ \cup\ \{\,0 \text{ with multiplicity }
\textstyle\sum_{j=1}^{r}(m_j-1)=|X|-r\,\}.
\]
Its maximum is $\max_j m_j=\max_{y\in Y}|f^{-1}(y)|$ (empty fibers contribute
$0$ to the maximum on the right and nothing to the block decomposition, so
including them changes nothing when some fiber is nonempty, which holds
because $X$ is nonempty and $f$ total). Hence
$\lVert A_f\rVert^2=\max_y|f^{-1}(y)|$, and taking square roots gives the
norm itself. The empty-domain case: $H_X$ is the zero space, every $A_f u$
with $u$ in the zero space is $0$, $\lVert A_f\rVert=0$ by convention, and
the fiber maximum over an empty source is $0$ under the stated reading.
\end{proof}

\begin{remark}[Full singular-value data]\label{ch07:rem:singular}
The proof delivers more than the norm: the complete singular-value multiset
of $A_f$ is $\{\sqrt{m_1},\dots,\sqrt{m_r}\}$ together with $|X|-r$ zeros,
since the singular values are the square roots of the Gram eigenvalues
\cite{horn}. Every spectral invariant of $A_f$ that depends only on singular
values is therefore an exact function of the fiber-size multiset. This is an
expository consequence of the proof, not a separately registered ledger
item; we state it because it explains \emph{why} the implementation exposes
\texttt{fiber\_sizes} as the primitive property from which
\texttt{max\_fiber\_size}, \texttt{operator\_norm\_squared}, \texttt{rank},
and the injectivity tests are all derived: the fiber-size tuple \emph{is}
the spectral data.
\end{remark}

\begin{example}[The theorem on the quartet]\label{ch07:ex:normquartet}
For $f$ (fibers $2,2,0$): $\lVert A_f\rVert^2=2$, Gram spectrum
$\{2,2,0,0\}$ --- and indeed $G_f$ in Example~\ref{ch07:ex:gramquartet} is
$J_2\oplus J_2$ already in block form, each block contributing eigenvalues
$\{2,0\}$. For $g$ (fibers $1,0,1$): $\lVert A_g\rVert^2=1$; $G_g=I_2$ with
spectrum $\{1,1\}$; $g$ is an isometry, as
Corollary~\ref{ch07:cor:isometry} will confirm from injectivity. For $\pi$:
$\lVert A_\pi\rVert^2=1$; permutation matrices are unitary. For $c_0$
(fibers $2,0$): $\lVert A_{c_0}\rVert^2=2$, so
$\lVert A_{c_0}\rVert=\sqrt{2}>1$ --- the noncontractivity that
Section~\ref{ch07:sec:counterexample} turns into a permanent fixture.
\end{example}

\subsection{A second proof, avoiding the spectral theorem}\label{ch07:sec:secondproof}

The block-Gram proof above is the registered one, and its route through the
spectral machinery is what makes the corollaries fall out for free. But the
norm theorem itself admits a shorter, wholly elementary proof --- two
inequalities, one application of Cauchy--Schwarz, one witness vector ---
and recording it serves two purposes: it double-derives a load-bearing
result, and it identifies the natural candidate route for the open Lean
formalization, since it needs neither eigenvalues, nor permutation
similarity, nor direct sums.

\begin{proof}[Second proof of Theorem~\textup{\ref{ch07:thm:norm}}]
Write $M=\max_{y}|f^{-1}(y)|$ and let $u=\sum_x c_xe_x$ be arbitrary.
Grouping the image coordinates by fiber,
\[
\lVert A_fu\rVert^2
=\sum_{y\in Y}\Bigl|\sum_{x\in f^{-1}(y)}c_x\Bigr|^2 .
\]
For each $y$ with fiber size $m_y$, Cauchy--Schwarz applied to the
$m_y$-term sum against the all-ones weight gives
$\bigl|\sum_{x\in f^{-1}(y)}c_x\bigr|^2
\le m_y\sum_{x\in f^{-1}(y)}|c_x|^2$.
Summing over $y$ and using $m_y\le M$ together with the disjointness of
fibers,
\[
\lVert A_fu\rVert^2
\le \sum_y m_y\!\!\sum_{x\in f^{-1}(y)}\!\!|c_x|^2
\le M\sum_{x}|c_x|^2 = M\,\lVert u\rVert^2 ,
\]
so $\lVert A_f\rVert^2\le M$. For the reverse inequality, pick $y^\ast$
with $|f^{-1}(y^\ast)|=M$ (possible since $X$ is nonempty and $f$ total)
and let $u^\ast=\tfrac{1}{\sqrt M}\sum_{x\in f^{-1}(y^\ast)}e_x$, a unit
vector. Then $A_fu^\ast=\tfrac{M}{\sqrt M}e_{y^\ast}=\sqrt M\,e_{y^\ast}$,
of norm $\sqrt M$. Hence $\lVert A_f\rVert^2\ge M$, and equality holds.
\end{proof}

Note where the equality conditions live: Cauchy--Schwarz is tight exactly
on constant coefficient patterns within a fiber, which is why the extremal
vector is the uniform fiber indicator --- the same vector the spectral
proof produces as the top eigenvector. The two proofs are not merely
consistent; they locate the same extremal geometry by different means. A
third, weaker route deserves a warning label: bounding
$\lambda_{\max}(G_f)$ by the maximum row sum of $G_f$ (a Gershgorin-type
bound) also yields $\le M$, since each row of $G_f$ contains exactly
$m_{f(x)}$ ones --- but a row-sum bound alone proves only the upper half,
and quoting it as the theorem silently drops the attainment argument. The
counterexample culture of this kernel applies to proofs too: half a proof
with a confident name attached is how false lemmas enter ledgers.

\begin{example}[The spectral decomposition, made explicit]\label{ch07:ex:eigenvectors}
For the running arrow $f$ it is worth writing the eigenvectors down rather
than merely counting them, because they carry the fiber combinatorics on
their sleeves. The Gram $G_f=J_2\oplus J_2$ (blocks on $\{p,q\}$ and
$\{r,s\}$) has the orthogonal eigenbasis
\[
v_1=\tfrac{1}{\sqrt2}(1,1,0,0)^T,\quad
v_2=\tfrac{1}{\sqrt2}(0,0,1,1)^T,\quad
v_3=\tfrac{1}{\sqrt2}(1,-1,0,0)^T,\quad
v_4=\tfrac{1}{\sqrt2}(0,0,1,-1)^T,
\]
with $G_fv_1=2v_1$, $G_fv_2=2v_2$, $G_fv_3=0$, $G_fv_4=0$. The
eigenvalue-$2$ vectors are the \emph{normalized fiber indicators}: uniform
mass across a fiber is what the pushforward amplifies maximally
($A_fv_1=\tfrac{1}{\sqrt2}(e_a+e_a)=\tfrac{2}{\sqrt2}\,e_a
=\sqrt{2}\,e_a$, of norm $\sqrt2$). The
kernel vectors are the \emph{within-fiber differences}: mass patterns that
cancel under merging ($A_fv_3=\tfrac{1}{\sqrt2}(e_a-e_a)=0$). The
$|X|-r=2$ dimensions of cancellation are exactly the information the
non-injective $f$ destroys, and the singular values $\sqrt2,\sqrt2$ on the
surviving directions are exactly the amplification the norm theorem
reports. Nothing about this picture is specific to $f$: for any arrow, the
normalized indicator of each nonempty fiber is a singular vector with
singular value the square root of the fiber size, and the orthogonal
complement inside each fiber block is killed.
\end{example}

The implementation computes \texttt{operator\_norm\_squared} as
\texttt{max\_fiber\_size} --- an exact integer, no linear algebra performed
at runtime --- and this identification is \emph{justified} by
Theorem~\ref{ch07:thm:norm}; the executable suite's structural tests check
the integer pipeline (fibers, maximum, squared norm) and cross-check derived
quantities on bounded exhaustive families, but the identification of
``maximum fiber size'' with ``operator norm squared'' rests on the informal
proof above until the open universal formalization closes it.

\section{Corollaries with proofs}\label{ch07:sec:corollaries}

Three consequences are implemented as first-class properties of
\texttt{DeterministicLift} and registered as implemented consequences of the
norm analysis. We prove each from the structure already established.

\begin{corollary}[Trace of the input Gram]\label{ch07:cor:trace}
$\trace(A_f^*A_f)=|X|$.
\end{corollary}

\begin{proof}
Two proofs, agreeing as they must. \emph{Diagonal count:} by
Proposition~\ref{ch07:prop:gramentries} every diagonal entry of $G_f$ is $1$
(each input collides with itself), and there are $|X|$ of them.
\emph{Spectral count:} by the proof of Theorem~\ref{ch07:thm:norm} the
eigenvalues of $G_f$ are the fiber sizes $m_1,\dots,m_r$ and zeros, and the
trace is their sum $\sum_j m_j$, which is $|X|$ because $f$ is total: the
nonempty fibers partition the source.
\end{proof}

The implementation property \texttt{trace\_input\_gram} returns
\texttt{len(arrow.source)} directly --- the corollary is what entitles it to
do so without forming any matrix.

\begin{corollary}[Rank]\label{ch07:cor:rank}
$\rk(A_f)=|\im(f)|$, the number of nonempty fibers.
\end{corollary}

\begin{proof}
The columns of $A_f$ are the vectors $e_{f(x)}$ for $x$ in the source
carrier; the set of distinct columns is $\{e_y : y\in \im f\}$, which is
orthonormal, hence linearly independent, and spans the column space. So the
column space has dimension $|\im f|$. (Spectrally: the rank of $A_f$ equals
the number of nonzero singular values, which by
Remark~\ref{ch07:rem:singular} is the number $r$ of nonempty fibers, and
$r=|\im f|$.)
\end{proof}

The implementation counts nonzero entries of \texttt{fiber\_sizes};
Corollary~\ref{ch07:cor:rank} is the license for that shortcut.

\begin{corollary}[Isometry criterion]\label{ch07:cor:isometry}
$A_f$ is an isometry ($A_f^*A_f=I$) if and only if $f$ is injective.
Consequently $\lVert A_f\rVert=1$ iff $f$ is injective (for nonempty $X$),
and $\lVert A_f\rVert>1$ iff $f$ has a collision.
\end{corollary}

\begin{proof}
By Proposition~\ref{ch07:prop:gramentries}, $G_f=I$ iff all off-diagonal
entries vanish (diagonal entries are automatically $1$), iff no two distinct
inputs share a fiber, iff $f$ is injective. For the norm statement: if $f$ is
injective all fibers have size $\le 1$ and some fiber has size exactly $1$
(nonempty source, total map), so the fiber maximum is $1$ and
$\lVert A_f\rVert=1$ by Theorem~\ref{ch07:thm:norm}; if $f$ is not injective
some fiber has size $\ge 2$ and $\lVert A_f\rVert\ge\sqrt{2}$.
\end{proof}

\begin{corollary}[No deterministic lift is a strict contraction]\label{ch07:cor:nocontraction}
For nonempty $X$, $\lVert A_f\rVert\ge 1$ always, with equality iff $f$ is
injective. There is no arrow whose lift has operator norm strictly below
$1$.
\end{corollary}

\begin{proof}
Immediate from Corollary~\ref{ch07:cor:isometry}: the fiber maximum of a
total map on a nonempty source is at least $1$.
\end{proof}

\begin{corollary}[Partial-isometry criterion]\label{ch07:cor:partialisometry}
$A_f$ is a partial isometry ($A_fA_f^*A_f=A_f$) if and only if $f$ is
injective --- i.e., for total deterministic lifts, partial isometry and
isometry coincide.
\end{corollary}

\begin{proof}
Compute on basis vectors, using
Remark~\ref{ch07:rem:adjointpullback} and
Proposition~\ref{ch07:prop:outputgram}:
\[
A_fA_f^*A_f\,e_x
= A_fA_f^*\,e_{f(x)}
= A_f\!\!\sum_{x'\,:\,f(x')=f(x)}\!\! e_{x'}
= \bigl|f^{-1}(f(x))\bigr|\; e_{f(x)} .
\]
This equals $A_fe_x=e_{f(x)}$ for all $x$ iff every \emph{inhabited} fiber
has size exactly $1$, iff $f$ is injective. (Empty fibers impose no
condition; but a total map's inhabited fibers cover the source.) The
implementation's \texttt{is\_partial\_isometry} therefore returns the same
Boolean as \texttt{is\_injective}, and its docstring says exactly this:
``for a total standard lift, equivalently whether $f$ is injective.''
\end{proof}

\begin{corollary}[Frobenius comparison]\label{ch07:cor:frobenius}
The Frobenius (Hilbert--Schmidt) norm of a lift satisfies
$\lVert A_f\rVert_F^2=|X|$, and the chain
\[
\lVert A_f\rVert \;\le\; \lVert A_f\rVert_F \;\le\;
\sqrt{\rk(A_f)}\;\lVert A_f\rVert
\]
holds with: left equality iff $f$ has exactly one nonempty fiber (i.e.\ $f$
is constant, or $X$ is empty); right equality iff all nonempty fibers have
equal size.
\end{corollary}

\begin{proof}
$\lVert A_f\rVert_F^2$ is the sum of squared entries; the entries are $0$
and $1$ with exactly one $1$ per column, so the sum is the number of
columns, $|X|$. (Equivalently
$\lVert A_f\rVert_F^2=\trace(A_f^*A_f)=|X|$ by
Corollary~\ref{ch07:cor:trace}.) In fiber terms the chain reads
\[
\max_j m_j \;\le\; \sum_j m_j \;\le\; r\,\max_j m_j
\]
after squaring, with $m_1,\dots,m_r$ the nonempty fiber sizes: the middle
quantity is $|X|$ and the outer ones are $\lVert A_f\rVert^2$ and
$\rk(A_f)\lVert A_f\rVert^2$ by Theorem~\ref{ch07:thm:norm} and
Corollary~\ref{ch07:cor:rank}. A maximum equals the sum of positive terms
iff there is one term; the sum equals $r$ times the maximum iff all terms
equal the maximum.
\end{proof}

The corollary is a useful sanity clamp in exact arithmetic: the integer
identities $\lVert A_f\rVert_F^2=|X|$ and
$\lVert A_f\rVert^2=\max_j m_j$ can be cross-checked against each other and
against the rank without any floating point, and the equality cases are
recognizable combinatorial conditions (constant map; equal-fiber map)
rather than numerical coincidences.

\begin{remark}[Why Corollary~\ref{ch07:cor:partialisometry} is not vacuous]
For general matrices, ``partial isometry'' is strictly weaker than
``isometry'' (any orthogonal projector is a partial isometry but not an
isometry unless it is the identity). The corollary says the deterministic
lift class is too rigid to see the difference: totality forces every input
to occupy some fiber, so the only way $A_f^*A_f$ can be a projector is to be
the identity. This collapse is a genuine fact about the lift class, not a
definitional triviality, and the counterexample below shows the other side:
an idempotent lift that fails to be a partial isometry at all.
\end{remark}

\section{Composition, fibers, and submultiplicativity}\label{ch07:sec:composition}

The norm theorem concerns one arrow; this short section records how the
norm interacts with composition, because the interaction is where the
fiber picture earns its keep a second time. Everything here is an
expository consequence of Theorem~\ref{ch07:thm:norm} and elementary fiber
combinatorics; nothing is a separately registered ledger item.

\begin{proposition}[Fibers of a composite]\label{ch07:prop:fibercomp}
Let $f:X\to Y$ and $g:Y\to Z$ be composable typed arrows. For every $z$ in
the carrier of $Z$,
\[
(g\circ f)^{-1}(z)\;=\;\bigsqcup_{y\in g^{-1}(z)} f^{-1}(y)
\qquad\text{(a disjoint union), hence}\qquad
\bigl|(g\circ f)^{-1}(z)\bigr|=\sum_{y\in g^{-1}(z)}\bigl|f^{-1}(y)\bigr| .
\]
Consequently
\[
\max_z\bigl|(g\circ f)^{-1}(z)\bigr|
\;\le\;
\Bigl(\max_y\bigl|f^{-1}(y)\bigr|\Bigr)\cdot
\Bigl(\max_z\bigl|g^{-1}(z)\bigr|\Bigr).
\]
\end{proposition}

\begin{proof}
An input $x$ satisfies $g(f(x))=z$ iff $f(x)$ lies in $g^{-1}(z)$, which
partitions the composite fiber by the intermediate value $f(x)$; the union
is disjoint because $f$ is a function. For the inequality: each summand is
at most $\max_y|f^{-1}(y)|$, and the number of summands is
$|g^{-1}(z)|\le\max_z|g^{-1}(z)|$.
\end{proof}

\begin{corollary}[Submultiplicativity, with exact witnesses]\label{ch07:cor:submult}
$\lVert A_{g\circ f}\rVert\le\lVert A_g\rVert\,\lVert A_f\rVert$.
\end{corollary}

\begin{proof}
Combine Proposition~\ref{ch07:prop:fibercomp} with
Theorem~\ref{ch07:thm:norm} on all three arrows; alternatively, apply the
standard submultiplicativity of the operator norm \cite{horn} to
$A_{g\circ f}=A_gA_f$ (Theorem~\ref{ch07:thm:functor}). That the two routes
agree is itself a small consistency check between the combinatorial and the
operator-theoretic sides of the chapter.
\end{proof}

\begin{example}[Strictness]\label{ch07:ex:strict}
The inequality is frequently strict, and the fiber formula says exactly
when: strictness occurs unless some single $z$ simultaneously has a
maximum-size $g$-fiber \emph{all of whose} elements carry maximum-size
$f$-fibers. Take $X_3'=(\texttt{Tri},(t_0,t_1,t_2))$, $X_2'=(\texttt{Bit},(0,1))$,
$X_1=(\texttt{Pt},(\ast))$; let $f:X_3'\to X_2'$ send $t_0,t_1\mapsto 0$ and
$t_2\mapsto 1$ (max fiber $2$), and let $g:X_2'\to X_1$ be the unique map
(max fiber $2$). The composite is the unique map $X_3'\to X_1$ with one
fiber of size $3$, and indeed
\[
\sqrt3=\lVert A_{g\circ f}\rVert<\lVert A_g\rVert\,\lVert A_f\rVert=\sqrt2\cdot\sqrt2=2 :
\]
the bound $2\cdot 2=4$ on the squared norm overcounts because $g$'s big
fiber contains one small $f$-fiber ($|f^{-1}(1)|=1$), and the exact
composite fiber size is the \emph{sum} $2+1=3$, not the product. The sum
formula of Proposition~\ref{ch07:prop:fibercomp} is sharper than the
product bound in every mixed case, and the implementation can afford the
sharp value always: it simply computes the composite arrow and reads its
fibers, never estimating.
\end{example}

\begin{proposition}[Normal lifts are exactly permutation lifts]\label{ch07:prop:normal}
Let $f:X\to X$ be an endomap, so that $A_f$ is a square matrix and
normality ($A_f^*A_f=A_fA_f^*$) is meaningful. Then $A_f$ is normal iff
$f$ is a bijection, in which case $A_f$ is unitary.
\end{proposition}

\begin{proof}
By Propositions~\ref{ch07:prop:gramentries} and~\ref{ch07:prop:outputgram},
$A_f^*A_f$ has unit diagonal while $A_fA_f^*$ is diagonal with the fiber
sizes. If the two agree then $A_f^*A_f$ is a diagonal matrix with unit
diagonal, i.e.\ the identity, so $f$ is injective
(Corollary~\ref{ch07:cor:isometry}); an injective endomap of a finite
carrier is a bijection; and then $A_f$ is a permutation matrix, which is
unitary, hence normal. The converse inclusion (unitary implies normal) is
immediate.
\end{proof}

\begin{remark}[Spectral radius versus norm]\label{ch07:rem:spectralradius}
For any square matrix the spectral radius satisfies
$\rho(A)\le\lVert A\rVert$, with equality for normal matrices \cite{horn}.
The lift class realizes the gap concretely: the counterexample operator
$A_{c_0}$ of the next section is idempotent, so its eigenvalues lie in
$\{0,1\}$ and $\rho(A_{c_0})=1$, while
$\lVert A_{c_0}\rVert=\sqrt2$. By
Proposition~\ref{ch07:prop:normal} this gap is the \emph{generic} situation
for endomap lifts --- any collision destroys normality --- and it is one
more reason no iteration-convergence story attaches to lifts: the
``eigenvalues bounded by one'' fact that powers such stories controls
$\rho$, not $\lVert\cdot\rVert$, and controls nothing about transient
amplification, which is exactly what the norm measures and exactly what
fiber collisions produce.
\end{remark}

\section{The permanent counterexample}\label{ch07:sec:counterexample}

Ledger record \textsc{thm-lift-occupation-counterexample}. This is the single
most quoted fixture in the kernel's test suite, and the Lean seed
kernel-checks this exact two-state expansion and its noncontractivity. Its
role is prophylactic: it permanently blocks two general claims that are
individually plausible, jointly tempting, and both false.

\subsection{The fixture, verified property by property}

Take the running arrow (R4): the constant map $c_0$ on the two-element
object $X_2'=(\texttt{Bit},(0,1))$, $c_0(0)=c_0(1)=0$. Write $A=A_{c_0}$:
\[
A=\begin{pmatrix}1&1\\0&0\end{pmatrix},
\qquad
A^{*}=\begin{pmatrix}1&0\\1&0\end{pmatrix}.
\]
We verify each claimed property by explicit computation.

\emph{(a) $A$ is idempotent: $A^2=A$.}
\[
A^2=\begin{pmatrix}1&1\\0&0\end{pmatrix}\begin{pmatrix}1&1\\0&0\end{pmatrix}
=\begin{pmatrix}1\cdot 1+1\cdot 0 & 1\cdot 1+1\cdot 0\\ 0&0\end{pmatrix}
=\begin{pmatrix}1&1\\0&0\end{pmatrix}=A .
\]
Categorically this mirrors $c_0\circ c_0=c_0$: a constant map composed with
itself is itself, and by functoriality
(Theorem~\ref{ch07:thm:functor} below) $A^2=A_{c_0\circ c_0}=A_{c_0}=A$.
The implementation's \texttt{is\_idempotent} checks the combinatorial form
of this condition --- the arrow is an endomap and every image point is a
fixed point (\texttt{mapping[image] == image} for each image) --- which is
exactly $f\circ f=f$.

\emph{(b) $\lVert A\rVert=\sqrt{2}$.} By Theorem~\ref{ch07:thm:norm}: the
fiber sizes are $(2,0)$, the maximum is $2$, the norm is $\sqrt 2$.
Directly, without the theorem: $A^*A=J_2$ (computed below), and
$J_2$ has largest eigenvalue $2$ by Lemma~\ref{ch07:lem:allones}, so
$\lVert A\rVert^2=2$ by Lemma~\ref{ch07:lem:gramreduction}. Still more
directly: the unit vector $u=\tfrac{1}{\sqrt2}(1,1)^T$ has
$Au=\tfrac{1}{\sqrt2}(2,0)^T$ with $\lVert Au\rVert=\sqrt2$, and no unit
vector can do better because
$\lVert A(u_1,u_2)^T\rVert=|u_1+u_2|\le\sqrt{2}\,\lVert u\rVert$ by
Cauchy--Schwarz. Three routes, one answer.

\emph{(c) $\trace(A^*A)=2$.}
\[
A^{*}A=\begin{pmatrix}1&0\\1&0\end{pmatrix}
\begin{pmatrix}1&1\\0&0\end{pmatrix}
=\begin{pmatrix}1&1\\1&1\end{pmatrix}=J_2,
\qquad \trace(J_2)=2=|X_2'| ,
\]
confirming Corollary~\ref{ch07:cor:trace} on the fixture.

\emph{(d) $\rk A=1$.} The two columns of $A$ are equal (both $e_0$), so the
column space is the line through $e_0$: rank $1=|\im c_0|$, confirming
Corollary~\ref{ch07:cor:rank}.

\emph{(e) $A$ is not self-adjoint, hence not an orthogonal projector.}
$A\neq A^{*}$ by inspection: the $(1,2)$ entry of $A$ is $1$ while the
$(1,2)$ entry of $A^{*}$ is $0$. Since an orthogonal projector must satisfy
\emph{both} $Q^2=Q$ and $Q^{*}=Q$ (Definition~\ref{ch07:def:projpi}), and
$A$ satisfies the first but not the second, $A$ is idempotent without being
a projector. (The implementation's \texttt{is\_orthogonal\_projector}
returns true only for identity endomaps, which
Proposition~\ref{ch07:prop:coincide} and this computation jointly justify:
a total deterministic lift on one declared carrier that is both idempotent
and self-adjoint must have symmetric $0/1$ matrix with one $1$ per column
fixing all image points --- forcing the identity.)

\emph{(f) $A$ is not a partial isometry.}
\[
AA^{*}=\begin{pmatrix}1&1\\0&0\end{pmatrix}\begin{pmatrix}1&0\\1&0\end{pmatrix}
=\begin{pmatrix}2&0\\0&0\end{pmatrix},
\qquad
AA^{*}A=\begin{pmatrix}2&0\\0&0\end{pmatrix}
\begin{pmatrix}1&1\\0&0\end{pmatrix}
=\begin{pmatrix}2&2\\0&0\end{pmatrix}=2A\neq A .
\]
So $AA^*A=2A$, off by exactly the fiber size $2$, as
Corollary~\ref{ch07:cor:partialisometry}'s basis computation predicts.
Equivalently, $A^*A=J_2$ is not a projector ($J_2^2=2J_2\neq J_2$).

For completeness, the three derived operators of
Section~\ref{ch07:sec:derived} on the fixture: input Gram $G_{c_0}=J_2$;
output Gram $\widehat{G}_{c_0}=\operatorname{diag}(2,0)$; support projector
$S_{c_0}=\operatorname{diag}(1,0)$. All three are distinct, and only the
third is a projector --- the fixture separates the trio maximally.

\subsection{The two false general claims the fixture blocks}

\begin{boundary}[Blocked claim 1: ``deterministic lifts are contractive'']\label{ch07:bd:contractive}
The claim would assert $\lVert A_f\rVert\le 1$ for every arrow $f$. It is
plausible because the matrix of $A_f$ is a $0/1$ matrix with unit column
sums, and column-stochastic matrices are famously ``non-expanding'' in the
folklore sense. It is false: $\lVert A_{c_0}\rVert=\sqrt2$, verified three
ways in (b) above. The folklore intuition belongs to the $\ell^1$ geometry
of probability vectors, not to the $\ell^2$ geometry of Hilbert space
(Section~\ref{ch07:sec:notstochastic} dissects this precisely). Any
downstream argument that needs $\lVert A_f\rVert\le 1$ must therefore prove
injectivity of $f$ first (Corollary~\ref{ch07:cor:isometry}); by
Corollary~\ref{ch07:cor:nocontraction} nothing better than equality with
$1$ is ever available.
\end{boundary}

\begin{boundary}[Blocked claim 2: ``idempotent lifts are projectors'']\label{ch07:bd:projector}
The claim would assert that $A_f^2=A_f$ implies $A_f$ is an orthogonal
projector. It is plausible because in many familiar operator classes
(self-adjoint operators, normal operators) idempotence does imply
projectorhood. It is false for lifts: $A_{c_0}$ is idempotent, not
self-adjoint, not a projector, and not even a partial isometry, by (a),
(e), (f) above. Structurally: idempotent \emph{arrows} $f\circ f=f$
abound (every constant map, every retraction onto a fixed-point set), and
their lifts are idempotent by functoriality, but self-adjointness of a lift
is an extremely rigid condition (symmetric $0/1$ matrix, one $1$ per
column) that combinatorially forces the identity among endomaps. The
implementation encodes the asymmetry as two distinct predicates ---
\texttt{is\_idempotent} and \texttt{is\_orthogonal\_projector} --- whose
truth values differ on this fixture, permanently, in the regression suite.
\end{boundary}

\begin{remark}[Why ``permanent fixture'' is the right category]
The fixture is registered in the theorem ledger
(\textsc{thm-lift-occupation-counterexample}) and kernel-checked in Lean,
which is an unusual amount of ceremony for a $2\times 2$ computation. The
ceremony is the point. The two blocked claims are exactly the kind of
plausible-sounding lemma that gets silently assumed in the middle of a long
argument about ``projection layers'' or ``contraction mappings'' in
cognitive-architecture prose. A counterexample that is registered, hashed,
tested on every run, and formally verified cannot be argued away, forgotten,
or weakened in a refactor without tripping an alarm. It converts a
mathematical fact into an engineering invariant.
\end{remark}

\section{Functoriality}\label{ch07:sec:functor}

\subsection{The theorem}

\begin{theorem}[Lift functor laws]\label{ch07:thm:functor}
For every object $X$, $A_{\id_X}=I_{H_X}$; and whenever the composite
$g\circ f$ is defined in the typed category (declared target of $f$ equals
declared source of $g$, as objects with name and carrier),
\[
A_{g\circ f}=A_g\,A_f .
\]
Consequently $A_{(-)}$ is a functor from the typed finite-function category
to the category of finite-dimensional Hilbert spaces and linear maps, in the
sense of~\cite{maclanecwm}: it assigns $H_X$ to $X$ and $A_f$ to $f$,
preserving identities and composition.
\end{theorem}

\begin{proof}
On basis vectors: $A_{\id_X}e_x=e_{\id_X(x)}=e_x=I_{H_X}e_x$ for every
carrier element $x$, and
\[
A_gA_f\,e_x = A_g\,e_{f(x)} = e_{g(f(x))} = e_{(g\circ f)(x)}
= A_{g\circ f}\,e_x .
\]
Two linear maps that agree on a basis agree everywhere, by linearity; and
both sides are linear by construction. The typing side condition matters:
$A_ge_{f(x)}$ is only meaningful because $f(x)$ is a standard index of the
object that is \emph{simultaneously} the declared target of $f$ and the
declared source of $g$ --- the same object, with the same name and the same
ordered carrier, hence the same basis indexed the same way. If the kernel
permitted cardinality-only coercion between distinct objects of equal
carrier size, the displayed computation would silently identify two
different bases and the middle equality would be unjustified. Functor
preservation of longer composites follows by induction from the binary case
together with associativity (Chapter~\ref{ch04:main}).
\end{proof}

\begin{example}[Functoriality on the running examples]\label{ch07:ex:functorexample}
Let $h:X_2\to X_4$ send $u\mapsto q$ and $v\mapsto s$ (injective), and
recall $f:X_4\to Y_3$ from (R1). The composite $f\circ h:X_2\to Y_3$ sends
$u\mapsto f(q)=a$ and $v\mapsto f(s)=b$. The three matrices:
\[
A_h=\begin{pmatrix}0&0\\1&0\\0&0\\0&1\end{pmatrix},\qquad
A_f=\begin{pmatrix}1&1&0&0\\0&0&1&1\\0&0&0&0\end{pmatrix},\qquad
A_{f\circ h}=\begin{pmatrix}1&0\\0&1\\0&0\end{pmatrix}.
\]
Multiplying: row $a$ of $A_f$ dotted with column $u$ of $A_h$ is
$1{\cdot}0+1{\cdot}1+0{\cdot}0+0{\cdot}0=1$; with column $v$:
$1{\cdot}0+1{\cdot}0+0{\cdot}0+0{\cdot}1=0$. Row $b$: with $u$, $0$; with
$v$, $1$. Row $c$: zeros. So
$A_fA_h=\left(\begin{smallmatrix}1&0\\0&1\\0&0\end{smallmatrix}\right)
=A_{f\circ h}$, an exact integer identity --- the kind of identity the
executable evidence checks wholesale.
\end{example}

\begin{remark}[Powers of endomaps]\label{ch07:rem:powers}
For an endomap $f:X\to X$, functoriality gives $A_{f^n}=A_f^{\,n}$ for all
$n\ge 0$ by induction ($A_{f^0}=A_{\id_X}=I$). In particular the
factorial-power idempotence theorem of Chapter~\ref{ch05:main} transports:
$A_f^{\,|X|!}$ is an idempotent matrix, being the lift of the idempotent
endomap $f^{|X|!}$. Note carefully what does \emph{not} transport: by
Boundary~\ref{ch07:bd:projector}, that idempotent matrix is generally not a
projector, and no convergence, mixing, or dynamical statement accompanies
it. Semigroupoid paths, having no power operation by type
(Chapter 6), induce no power identity here either: the lift of
a path composite is the product of the lifts, and that is all.
\end{remark}

\begin{remark}[Faithful, not full]\label{ch07:rem:faithfulnotfull}
The lift functor is \emph{faithful} (Proposition~\ref{ch07:prop:matrix}:
distinct parallel arrows have distinct matrices) but very far from
\emph{full}: almost no linear map $H_X\to H_Y$ is a lift. Any matrix with an
entry outside $\{0,1\}$ (say $2I$, or the averaging map
$\tfrac12 J_2$), and any $0/1$ matrix with a zero column or a two-$1$
column, lies outside the image. Fullness failure is a feature, not a
defect: it is the formal statement that the linear world strictly extends
the functional world, and it marks exactly where extension proposals ---
stochastic columns, superposed columns --- would leave the licensed kernel
and enter open research territory (Section~\ref{ch07:sec:notclaims}). In
the terminology of~\cite{maclanecwm}, the kernel uses the lift as a
faithful comparison functor and never claims an equivalence of categories;
no downstream argument may pull a linear-algebraic construction back along
$A_{(-)}$ unless the construction demonstrably lands in the
one-$1$-per-column image.
\end{remark}

\subsection{Executable evidence: the 1{,}678 pairs}

The functor laws are the most heavily instrumented claim in the lift module.
The executable suite verifies them by \emph{exact matrix equality} --- integer
entries, tuple-of-tuples comparison, no floating point anywhere --- for every
valid typed-map pair on carriers through size three: $1{,}678$ composition
pairs, including pairs involving empty carriers. Concretely, for each such
pair $(f,g)$ with matching intermediate object, the suite forms the integer
matrix product of \texttt{matrix()} outputs and compares it entrywise with
the \texttt{matrix()} of the composed arrow; identity lifts are compared
against exact identity matrices.

\begin{boundary}[What the 1{,}678 pairs do and do not establish]\label{ch07:bd:pairs}
Passing this family establishes the functor laws \emph{on the covered finite
family, for the running implementation, on the run in question}. It is
bounded exhaustive testing in the small-scope tradition~\cite{jackson}: within
the size bound, no case is sampled away. It is not a proof of the universal
statement, not proof-assistant verification of the Python code, and it
closes no ledger obligation. The universal statement's formal status is
recorded next.
\end{boundary}

\section{Formal verification scope: verified versus open}\label{ch07:sec:lean}

The monograph's policy is to state Lean scope exactly, per claim, with the
same care used for the claims themselves. The formal seed is a pinned,
isolated Lean~4.32.1 toolchain~\cite{lean4} whose verifier hardcodes the
sole allowed standard axiom (\texttt{propext}), rejects project axiom
declarations and \texttt{sorry}, checks source hashes, and audits exact
theorem dependency reports (Chapter 11 describes the harness).
For the material of this chapter:

\emph{Kernel-checked (verified scope).}
\begin{itemize}
\item The identity lift, \emph{entrywise}: the matrix of $A_{\id_X}$ equals
  the identity matrix entry by entry, and $A_{\id_X}$ fixes every basis
  vector.
\item Typed composition \emph{on basis vectors}: for composable $f,g$ in the
  Lean model, $A_gA_f$ and $A_{g\circ f}$ agree on every basis vector.
\item Relational-product \emph{support composition}: the $0/1$ support
  pattern of a composite behaves as the relational product of the supports.
\item The exact two-state expansion of the permanent counterexample
  (Section~\ref{ch07:sec:counterexample}) and its noncontractivity.
\end{itemize}

\emph{Open (not verified).}
\begin{itemize}
\item The universal norm theorem (Theorem~\ref{ch07:thm:norm}); obligation
  PO-TKS8-CORE-05 remains \textsc{informal proof present / universal
  formalization open}.
\item The generic matrix bridge: the connection, in full generality, between
  the basis-vector definition of the lift and the numerical matrix, Gram,
  trace, and rank definitions --- including finite-sum numerical matrix
  multiplication.
\item Arbitrary-vector linearity: the extension of the verified basis-vector
  identities to all vectors of the space, in the Lean model.
\end{itemize}

\begin{remark}[On the shape of the open norm formalization]\label{ch07:rem:leanroute}
It is worth recording, as an engineering observation rather than a
commitment, why the universal norm theorem is the hard item on the open
list and what the path of least resistance looks like. The registered
block-Gram proof consumes three pieces of nontrivial formal
infrastructure: basis permutation as a similarity, spectra of direct sums,
and the spectral theorem for Hermitian matrices --- each routine on paper,
each a substantial formal dependency. The second proof of
Section~\ref{ch07:sec:secondproof} consumes only finite sums,
Cauchy--Schwarz, and one witness computation, all of which sit close to the
finite-sum matrix bridge that is already an open item for the functor laws.
A formalization that lands the bridge would therefore likely land the norm
theorem by the two-inequality route in the same effort, while the spectral
route would remain the source of the corollaries' \emph{spectral} phrasing
only. None of this closes anything: PO-TKS8-CORE-05 remains open, and this
remark is a map of the terrain, not a claim about it.
\end{remark}

Obligation PO-TKS8-CORE-11 (lift preserves identities and composition)
therefore carries the compound status \textsc{partial formal verification
present / finite-sum matrix bridge open}. Two remarks keep this scope honest.
First, ``basis vectors plus linearity'' is a complete informal argument
(Theorem~\ref{ch07:thm:functor}) but the linearity step is exactly what is
not yet formalized, so the Lean artifact verifies a strictly weaker
statement than the theorem. Second, the equivalence between the Lean model
and the Python implementation is itself an open obligation; a kernel-checked
Lean lemma constrains the Lean model, not the running code.

\section{What the lift is not}\label{ch07:sec:notclaims}

The lift's matrix shape invites three importations from neighboring theories
--- Markov chains, quantum information, and dynamical systems. All three are
declined, explicitly and permanently. This section explains each refusal
with enough mathematical detail that the boundary can be checked rather than
merely asserted.

\subsection{Not stochastic matrices}\label{ch07:sec:notstochastic}

Every lift matrix has nonnegative entries and unit column sums, so as a
\emph{matrix} it satisfies the definition of a column-stochastic matrix ---
in fact it is the extreme, deterministic kind, in which each column is a
point mass. The refusal is therefore not about the matrix class; it is about
which \emph{theory} of that matrix class is in force, and three separations
do the work.

\emph{Different norm, different geometry.} Stochastic-matrix theory lives on
the probability simplex with the $\ell^1$ norm, where every column-stochastic
matrix $M$ is nonexpanding on nonnegative vectors:
$\lVert Mv\rVert_1=\sum_y|\sum_x M_{yx}v_x|=\sum_x v_x\sum_y M_{yx}
=\lVert v\rVert_1$ for $v\ge 0$ --- mass is conserved. The lift theory of
this chapter lives on Hilbert space with the $\ell^2$ norm, where
Theorem~\ref{ch07:thm:norm} gives $\lVert A_f\rVert_2=\sqrt{\max_y m_y}$,
which exceeds $1$ the moment any fiber has two elements. Both statements are
true of the same matrix; they answer different questions. The blocked claim
of Boundary~\ref{ch07:bd:contractive} is precisely the illegitimate
transport of the first statement's conclusion into the second statement's
setting.

\emph{Different morphism class.} The deterministic lifts form a vanishingly
thin subclass of stochastic matrices: a column such as
$(\tfrac12,\tfrac12,0)^T$ is column-stochastic but is not $e_{f(x)}$ for any
$f$, so no genuinely randomizing kernel is a lift. Conversely, the kernel's
category has no arrows whose lift would need fractional entries. The functor
$A_{(-)}$ lands in stochastic matrices but is not surjective onto them, and
no completion of the image to general stochastic matrices is defined,
admitted, or licensed here; ``stochastic kernels'' appear on the excluded
list of the governing specification and remain \textsc{open}.

\emph{No probabilistic semantics.} Even where the arithmetic coincides ---
$A_f$ applied to a probability vector does compute the pushforward
distribution of $f$ --- the kernel asserts no semantics of uncertainty,
belief, frequency, or sampling. The coordinates that \texttt{apply} acts on
are coordinates, full stop. Reading them as probabilities is a semantic
transfer, and semantic transfers are exactly what the Ring~4 discipline
gates behind separate review.

\subsection{Not quantum channels}\label{ch07:sec:notquantum}

A quantum channel is a completely positive, trace-preserving (CPTP) map on
density operators. There is a standard way to try to manufacture one from
$A_f$: the conjugation map $\Phi_f(\rho)=A_f\,\rho\,A_f^{*}$. The map
$\Phi_f$ is completely positive (it has the single Kraus operator $A_f$),
but it is trace-preserving only in the degenerate case:
\[
\trace\bigl(A_f\rho A_f^{*}\bigr)=\trace\bigl(\rho\,A_f^{*}A_f\bigr)
=\trace(\rho\,G_f),
\]
which equals $\trace(\rho)$ for all $\rho$ iff $G_f=I$ iff $f$ is injective
(Corollary~\ref{ch07:cor:isometry}). On the permanent fixture the failure is
quantitative: with $\rho=\tfrac12 J_2$ (the density matrix of the uniform
superposition state on two basis states),
\[
\trace(\rho\,G_{c_0})=\trace\bigl(\tfrac12 J_2\cdot J_2\bigr)
=\tfrac12\trace(J_2^2)=\tfrac12\trace(2J_2)=\tfrac12\cdot 4=2\neq 1 .
\]
A ``channel'' that doubles the trace of a state is no channel. One could
repair trace preservation by renormalizing or by adding Kraus operators, but
any such repair is a new mathematical object with new semantics, not the
lift, and constructing it is not undertaken here: quantum channels sit on
the specification's excluded list alongside stochastic kernels, and the
noncontractivity fixture exists partly to keep the ``lifts are like
isometries, isometries are like channels'' slide from ever starting. The
only honest statement is the narrow one: for \emph{injective} $f$, $A_f$ is
an isometry and $\Phi_f$ is a legitimate isometric embedding of state
spaces; for anything else, no channel reading exists without modification.

\subsection{No dynamics}\label{ch07:sec:nodynamics}

Nothing in this chapter has a time parameter. The lift assigns one linear
map to one arrow; it does not assert that anything \emph{evolves} by
repeated application of that map. Three specific non-claims:

\emph{No trajectory semantics.} Iterating $A_f$ for an endomap $f$ computes
lifts of iterates ($A_f^n=A_{f^n}$, Remark~\ref{ch07:rem:powers}); it does
not model a system's state at time $n$, because no system, state, or clock
has been defined. The factorial-power idempotence result is a statement
about the monoid structure of finite endomaps, not about equilibration.

\emph{No fixed-point or convergence claims.} Spectral radius arguments
(``iterate and converge to the dominant eigenspace'') require normality or
at least diagonalizability assumptions plus a semantics for convergence;
the fixture $A_{c_0}$ already shows lifts can be non-normal
($A^*A=J_2\neq \operatorname{diag}(2,0)=AA^{*}$), and no convergence
semantics is defined.

\emph{No cognitive-dynamics reading.} The governing sources are explicit
that mapping canonical Noetic or CAS constructs onto kernel arrows is an
open, unlicensed obligation (PO-TKS8-CORE-08). A fortiori, reading
$A_f$-iteration as ``cognitive state evolution'' is doubly out of scope:
it would require both the semantic mapping and a dynamical interpretation,
neither of which exists. The lift is bookkeeping --- an exact,
faithful, spectrally transparent ledger of which inputs an arrow merges and
which outputs it reaches --- and the entire content of this chapter is what
can be proved about that ledger.

\section{Summary and status table}\label{ch07:sec:summary}

The chapter's results, with their exact statuses:

\begin{center}
\begin{tabular}{@{}llp{5.2cm}@{}}
\toprule
Result & Where & Status \\
\midrule
Lift well-defined; matrix bijection & Prop.~\ref{ch07:prop:matrix} & definition-level; immediate \\
Gram entries; fiber partition & Prop.~\ref{ch07:prop:gramentries} & immediate computation \\
Output Gram diagonal, fiber sizes & Prop.~\ref{ch07:prop:outputgram} & immediate computation \\
Support projector onto $\im A_f$ & Prop.~\ref{ch07:prop:supportisproj} & immediate computation \\
Trio coincidence iff injective & Prop.~\ref{ch07:prop:coincide} & immediate computation \\
Exact norm formula & Thm.~\ref{ch07:thm:norm} & informal proof present; universal formalization open (PO-05) \\
Trace, rank, isometry corollaries & Cor.~\ref{ch07:cor:trace}--\ref{ch07:cor:isometry} & consequences of the norm analysis \\
No strict contraction; partial-isometry collapse & Cor.~\ref{ch07:cor:nocontraction}, \ref{ch07:cor:partialisometry} & consequences; expository registration \\
Permanent counterexample & Sec.~\ref{ch07:sec:counterexample} & finite computation; Lean kernel-checked (exact two-state expansion) \\
Functor laws & Thm.~\ref{ch07:thm:functor} & partial formal verification; matrix bridge and linearity open (PO-11); $1{,}678$ exact-matrix pairs \\
Non-claims (stochastic, quantum, dynamics) & Sec.~\ref{ch07:sec:notclaims} & permanent boundaries \\
\bottomrule
\end{tabular}
\end{center}

\begin{exercise}\label{ch07:ex:exercise1}
Let $h:X_5\to Y_2$ be the arrow from a five-element carrier to a two-element
carrier with fibers of sizes $4$ and $1$. Without forming any matrix, write
down: $\lVert A_h\rVert$; the full spectrum of $A_h^*A_h$ with
multiplicities; $\trace(A_h^*A_h)$; $\rk(A_h)$; the diagonal of
$A_hA_h^*$; the support projector; and the truth values of
\texttt{is\_injective}, \texttt{is\_isometry},
\texttt{is\_partial\_isometry}. Then check
Corollary~\ref{ch07:cor:frobenius}'s chain numerically:
$\sqrt{4}\le\sqrt{5}\le\sqrt{2}\cdot\sqrt{4}$, and identify which equality
case fails and why.
\end{exercise}

\begin{exercise}\label{ch07:ex:exercise2}
Prove that the matrix of $A_f$ is doubly stochastic (all row sums and all
column sums equal to $1$) iff $f$ is a bijection between carriers of equal
size, and conclude via Proposition~\ref{ch07:prop:normal} that within the
lift class ``doubly stochastic,'' ``unitary,'' ``normal,'' and
``permutation matrix'' all coincide --- a collapse of four generally
distinct matrix classes into one, forced by the one-$1$-per-column shape.
\end{exercise}

\begin{exercise}\label{ch07:ex:exercise3}
The counterexample $A_{c_0}$ satisfies $A_{c_0}A_{c_0}^*A_{c_0}=2A_{c_0}$.
Show that in general, for an \emph{idempotent} arrow $f=f\circ f$ whose
image points all have fiber size $k$, one has $A_fA_f^*A_f=kA_f$ on the
image summand, and use the constant map on an $m$-element carrier to
exhibit $\lVert A_f\rVert=\sqrt{m}$ growing without bound over idempotent
lifts. (Moral: not only are idempotent lifts not projectors; they are not
even uniformly bounded.)
\end{exercise}

Everything above is Ring~4, draft, not licensed, with independent specialist
review open (PO-TKS8-CORE-10). The next chapter turns from the linear shadow
of single arrows to the combinatorics of arrow \emph{words}: which
compositions of a declared generator set produce which denotations, and how
to name each denotation by a canonical representative that is stable,
replayable, and honest about its resource limits.
